This follows immediately from 1.9 and from the proof of 1.5.

\sgaStableTarget{sga6.e08.sec.2}\subsection*{2. Relative cotangent complex}

\sgaStableTarget{sga6.e08.para.2-1}\textbf{2.1.} We shall give here a definition of the cotangent complex which, without having all the generality possible (see Exp. 0 4.4), will suffice for the needs of the Riemann--Roch theorem; of course, the definition given in exp. 0 coincides in the case considered with the definition given here.

Let @@M0@@ be a smoothable morphism, that is, one admitting a factorization:
@@M1@@
where @@M2@@ is a smooth morphism, and @@M3@@ an immersion which may be supposed closed. Let @@M4@@ be the ideal of @@M5@@ in @@M6@@. The canonical derivation:
@@M7@@
defines a homomorphism:
@@M8@@
which vanishes on @@M9@@, and hence defines a homomorphism, still denoted @@M10@@:
\sgaStableTarget{sga6.e08.eq.2-1-1}@@M11@@
One checks easily that this homomorphism is @@M12@@-linear. One then defines a chain complex @@M13@@ by putting:
@@M14@@
and taking as homomorphism @@M15@@ the homomorphism (2.1.1).
